\documentclass[11pt]{letter}

\usepackage[margin=1in]{geometry}
\usepackage{hyperref}
\emergencystretch=2em

\signature{Vincenzo Barone}
\address{Vincenzo Barone\\
INSTM, National Interuniversity Consortium for Materials Science and Technology\\
via G. Giusti 9, 50121 Firenze, Italy\\
\href{mailto:vincebarone52@gmail.com}{vincebarone52@gmail.com}}

\begin{document}

\begin{letter}{Editor\\
The Journal of Chemical Physics}

\opening{Dear Editor,}

Please find enclosed the manuscript ``Algorithmic Construction and Exploitation
of Topology-Constrained Refinement Models for Semiexperimental Equilibrium
Structures,'' submitted for consideration as an Article in \emph{The Journal of
Chemical Physics}.

The manuscript introduces MORPHEUS, a framework that turns the formulation of
semiexperimental refinement models into a reproducible topology-, symmetry-,
and rank-audited procedure. Rather than assuming that the internal-coordinate
model has already been selected by expert intervention, MORPHEUS constructs the
active coordinate space from the molecular graph, applies symmetry and
homogeneous rank reduction, projects constraints and predicates consistently,
and propagates uncertainties back to chemically interpretable structural
parameters.

The work is, in my view, well aligned with the scope of \emph{The Journal of
Chemical Physics} because it addresses a general problem at the interface of
rotational spectroscopy, molecular structure determination, and quantum
chemistry: how to combine high-resolution experimental data, computed
vibrational and electronic corrections, topology, symmetry, and external
reference information in a controlled statistical model. The validation spans
acyclic, cyclic, planar, fused, bridged, and flexible molecules. A
large-molecule testosterone case study further shows that accurate quantum
chemical equilibrium structures, accelerated vibrational corrections, and
algorithmically generated topology-constrained refinement models can be chained
for systems containing about 50 atoms, where conventional isotope-rich
semiexperimental analyses are not currently practical.

The manuscript also emphasizes reproducibility. The reported models are not
defined by molecule-specific Z-matrix choices, but by explicit primitive
relations, parameter classes, predicates, and audit trails. The submitted
Supporting Information documents the benchmark inputs, diagnostics, and the
large-molecule class/predicate variants used in the paper.

This manuscript is original, has not been published previously, and is not
under consideration elsewhere. I confirm that there are no conflicts of
interest to declare.

\closing{Sincerely,}

\end{letter}

\end{document}
